%%============================================================================
%% CÓDIGO DE CONDUTA — Grupo de Usuários Linux em Estatística e Ciência de Dados
%%                     Universidade Federal do Paraná (UFPR)
%%============================================================================
%% Documento institucional — versão formal
%% Compilar com: pdflatex codigo-conduta.tex  (ou xelatex / lualatex)
%%============================================================================

\documentclass[
  12pt,          % tamanho base da fonte
  a4paper,       % formato de papel
  oneside,       % impressão frente única
  brazil         % idioma principal: português do Brasil
]{article}

%---------------------------------------------------------------------------
% Pacotes essenciais
%---------------------------------------------------------------------------
\usepackage[T1]{fontenc}            % codificação de fonte
\usepackage[utf8]{inputenc}         % entrada UTF-8 (pdflatex)
\usepackage{lmodern}                % fontes Latin Modern (vetoriais)
\usepackage{microtype}              % microtipografia (melhora legibilidade)

\usepackage[brazil]{babel}          % hifenização e nomes de seções em PT-BR

\usepackage[
  margin=2.5cm,
  top=3cm,
  bottom=3cm,
  headheight=25pt
]{geometry}                         % margens do documento

\usepackage{setspace}               % controle de espaçamento entre linhas
\usepackage{xcolor}                 % suporte a cores
\usepackage{titlesec}               % formatação de títulos de seções
\usepackage{titling}                % formatação do título do documento
\usepackage{fancyhdr}               % cabeçalhos e rodapés personalizados
\usepackage{enumitem}               % personalização de listas
\usepackage{lettrine}               % capitulares (letra inicial decorada)
\usepackage{graphicx}               % suporte a imagens (se necessário)
\usepackage{hyperref}               % hiperlinks e metadados do PDF
\usepackage{bookmark}               % bookmarks do PDF mais limpos

%---------------------------------------------------------------------------
% Cores institucionais
%---------------------------------------------------------------------------
\definecolor{gulecdBlue}{HTML}{1B3A5C}   % azul escuro institucional
\definecolor{gulecdAccent}{HTML}{2E86AB} % azul de destaque
\definecolor{gulecdGray}{HTML}{555555}   % cinza para texto secundário
\definecolor{gulecdLight}{HTML}{F5F7FA}  % fundo claro para caixas
\definecolor{gulecdRule}{HTML}{D0D5DD}   % cor de filetes e réguas

%---------------------------------------------------------------------------
% Metadados do PDF
%---------------------------------------------------------------------------
\hypersetup{
  pdftitle    = {Código de Conduta — GULECD/UFPR},
  pdfauthor   = {Grupo de Usuários Linux - Estatística e Ciência de Dados - UFPR},
  pdfsubject  = {Código de Conduta institucional do GULECD},
  pdfkeywords = {código de conduta, software livre, código aberto, ciência de dados, UFPR},
  colorlinks  = true,
  linkcolor   = gulecdAccent,
  urlcolor    = gulecdAccent,
  citecolor   = gulecdAccent
}

%---------------------------------------------------------------------------
% Configurações tipográficas
%---------------------------------------------------------------------------
\onehalfspacing                     % espaçamento entre linhas = 1,5
\setlength{\parindent}{1.5cm}       % recuo de parágrafo
\setlength{\parskip}{0.5em}         % espaço entre parágrafos

%---------------------------------------------------------------------------
% Formato do título do documento
%---------------------------------------------------------------------------
\pretitle{%
  \begin{center}
  \vspace*{1cm}
  \rule{\textwidth}{2pt}\vspace{0.3cm}\\
  \LARGE\bfseries\color{gulecdBlue}
}
\posttitle{%
  \\\vspace{0.3cm}
  \rule{\textwidth}{2pt}
  \end{center}
  \vspace{0.5cm}
}
\preauthor{\begin{center}\large\color{gulecdGray}}
\postauthor{\end{center}}
\predate{\begin{center}\small\color{gulecdGray}}
\postdate{\end{center}}

%---------------------------------------------------------------------------
% Formato das seções
%---------------------------------------------------------------------------
\titleformat{\section}
  [block]
  {\normalfont\Large\bfseries\color{gulecdBlue}}
  {\thesection.}{0.8em}{}
  [\vspace{0.2em}{\color{gulecdRule}\rule{\linewidth}{0.6pt}}]

\titleformat{\subsection}
  [block]
  {\normalfont\large\bfseries\color{gulecdAccent}}
  {\thesubsection.}{0.8em}{}

\titlespacing*{\section}
  {0pt}{2.5ex plus 1ex minus .2ex}{1.5ex plus .2ex}
\titlespacing*{\subsection}
  {0pt}{2ex plus 1ex minus .2ex}{1ex plus .2ex}

%---------------------------------------------------------------------------
% Cabeçalho e rodapé
%---------------------------------------------------------------------------
\pagestyle{fancy}
\fancyhf{}
\fancyhead[L]{\small\color{gulecdGray}\textsc{GULECD/UFPR}}
\fancyhead[R]{\small\color{gulecdGray}\textit{Código de Conduta}}
\fancyfoot[C]{\small\color{gulecdGray}---\ \thepage\ ---}
\renewcommand{\headrulewidth}{0.4pt}
\renewcommand{\headrule}{%
  \hbox to\headwidth{%
    \color{gulecdRule}\leaders\hrule height \headrulewidth\hfill}}

% Remove cabeçalho/rodapé da primeira página (será redefinido manualmente)
\fancypagestyle{plain}{%
  \fancyhf{}%
  \fancyfoot[C]{\small\color{gulecdGray}---\ \thepage\ ---}%
  \renewcommand{\headrulewidth}{0pt}%
}

%---------------------------------------------------------------------------
% Ambiente personalizado para listas do código de conduta
%---------------------------------------------------------------------------
\newlist{conduta}{itemize}{1}
\setlist[conduta]{
  label      = {\color{gulecdAccent}\textbullet},
  leftmargin = 1.5em,
  itemsep    = 0.3em,
  topsep     = 0.4em
}

%---------------------------------------------------------------------------
% Informações do documento
%---------------------------------------------------------------------------
\title{Código de Conduta\\[0.3em]
  \Large\mdseries\color{gulecdGray}%
  Grupo de Usuários Linux e \textit{Software} Livre\\Estatística e Ciência de Dados - UFPR}
%\author{Grupo de Usuários Linux e Software Livre\\Curso de Estatística e Ciência de Dados --- UFPR}
\date{\today}

%============================================================================
% INÍCIO DO DOCUMENTO
%============================================================================
\begin{document}

%--- Capa / Título principal -------------------------------------------------
\thispagestyle{empty}
\maketitle

\newpage

%--- Seção 1: Compromisso ---------------------------------------------------
\section{Compromisso}

\lettrine[lines=3, lhang=0.1, nindent=0em, slope=-0.5em]{N}{ós}, enquanto membros, colaboradores e organizadores, comprometemo-nos a 
garantir um ambiente de convivência acadêmica e técnica aberto, respeitoso e colaborativo. Nosso objetivo é promover uma comunidade onde o 
aprendizado, a troca de conhecimento e a produção científica ocorram de forma ética, inclusiva e tecnicamente rigorosa.

%--- Seção 2: Padrões de Comportamento --------------------------------------
\section{Padrões de Comportamento}

São considerados comportamentos desejáveis:

\begin{conduta}
  \item Uso de linguagem respeitosa, técnica e construtiva
  \item Abertura à crítica fundamentada e ao debate científico rigoroso
  \item Compartilhamento de conhecimento, código e recursos de forma colaborativa
  \item Reconhecimento de contribuições alheias
  \item Apoio a membros em processo de aprendizado
  \item Comunicação clara e baseada em evidências
  \item Participação ativa nas discussões, deliberações e votações do grupo de maneira honesta e aberta
\end{conduta}

%--- Seção 3: Comportamentos Inaceitáveis -----------------------------------
\section{Comportamentos Inaceitáveis}

Não serão tolerados:

\begin{conduta}
  \item Assédio, intimidação ou discriminação de qualquer natureza
  \item Ataques pessoais ou desqualificação não técnica
  \item Linguagem ofensiva, agressiva ou desrespeitosa
  \item Divulgação indevida de informações pessoais ou sensíveis
  \item Sabotagem de discussões técnicas ou projetos colaborativos
  \item Qualquer conduta que comprometa a integridade científica ou ética do grupo
\end{conduta}

%--- Seção 4: Escopo de Aplicação -------------------------------------------
\section{Escopo de Aplicação}

Este Código de Conduta aplica-se a todos os espaços do grupo, incluindo:

\begin{conduta}
  \item Grupo de WhatsApp e demais canais digitais
  \item Repositórios de código e plataformas colaborativas
  \item Reuniões, eventos, workshops e encontros presenciais
  \item Interações externas representando o grupo
\end{conduta}

%--- Seção 5: Responsabilidades da Moderação --------------------------------
\section{Responsabilidades da Moderação}

Os moderadores e coordenadores são responsáveis por:

\begin{conduta}
  \item Garantir a aplicação deste Código de Conduta
  \item Avaliar denúncias com imparcialidade e confidencialidade
  \item Tomar medidas proporcionais às infrações
  \item Promover um ambiente seguro e produtivo para todos os membros
\end{conduta}

%--- Seção 6: Aplicação e Medidas Corretivas --------------------------------
\section{Aplicação e Medidas Corretivas}

Em caso de violação, poderão ser adotadas, conforme a gravidade:

\begin{conduta}
  \item Advertência privada
  \item Solicitação de retratação
  \item Restrição temporária de participação
  \item Remoção do grupo
  \item Encaminhamento institucional, quando aplicável
\end{conduta}

%--- Seção 7: Denúncias -----------------------------------------------------
\section{Denúncias}

Situações que violem este Código de Conduta devem ser reportadas aos moderadores ou coordenadores do grupo. 
Todas as denúncias serão tratadas com seriedade, confidencialidade e diligência.

%--- Seção 8: Princípios Científicos e Éticos --------------------------------
\section{Princípios Científicos e Éticos}

Além das normas de convivência, espera-se que os membros:

\begin{conduta}
  \item Mantenham rigor científico em suas contribuições
  \item Evitem práticas como plágio, fabricação ou manipulação indevida de dados
  \item Respeitem normas éticas aplicáveis, incluindo a Lei Geral de Proteção de Dados (LGPD)
  \item Declarem conflitos de interesse relevantes
  \item Citem fontes e referências adequadas
  \item Promovam reprodutibilidade, transparência e o uso de \textit{software livre} e código aberto
  
\end{conduta}

%--- Seção 9: Alinhamento com Software Livre --------------------------------
\section{Alinhamento com Software Livre}

Este grupo orienta-se pelos princípios do \textit{software} livre e código aberto, incentivando:

\begin{conduta}
  \item Liberdade de uso, estudo, modificação e distribuição de \textit{software}
  \item Transparência técnica e auditabilidade
  \item Colaboração descentralizada e meritocrática
  \item Inclusão e diversidade na participação 
  \item Respeito às licenças e à propriedade intelectual 
  \item Privacidade e segurança como prioridades 
  \item Compartilhamento de conhecimento
  \item Promoção de software livre e código aberto como pilares de modelos econômicos sustentáveis
\end{conduta}

%--- Seção 10: Disposições Finais -------------------------------------------
\section{Atualizações deste Documento}

Este documento deve ser entendido como um instrumento vivo, passível de aprimoramento conforme a evolução do grupo, mudanças operacionais ou alterações legais. 
Sua aplicação deve sempre priorizar os princípios fundamentais de abertura, participação, eficiência e responsabilidade coletiva.

Novas versões deste documento deverão ser arquivadas juntamente com as versões anteriores no repositório \textit{documentos-referendados} e numeradas em acordo
 com o esquema padrão de \textit{releases} como XX.yy, onde XX refere-se a edições substanciais e yy à numeração em versões de ajuste. 
 
Também deverá ser mantida ativa a versão em edição no repositório \textit{documents-dev} a fim de prover o rastreamento de todas as mudanças entre as versões do documento,

Assim, determina-se que:
\begin{conduta}
	\item Este regulamento poderá ser revisado mediante discussão coletiva
	\item Situações não previstas serão resolvidas pela Coordenação vigente
	\item As atualizações serão comunicadas de forma transparente aos membros
\end{conduta}

\vspace{2cm}

%============================================================================
\end{document}
%============================================================================
